% Archive record 63, line 1.
\hypertarget{part-4-63-0001-iv-c-financial-factors-machine-learning-models}{%
\section{IV-C --- Financial Factors \& Machine Learning Models}\label{part-4-63-0001-iv-c-financial-factors-machine-learning-models}}

上一阶段 IV-B，我们证明了一个非常重要的数学关系：
\[R_{\mathrm{optimal}}^2
=
\frac{\mathrm{SNR}}{1+\mathrm{SNR}}.\]

这说明，在给定信息集合下，金融预测能力存在理论上限。

但在实际 Quantitative Research 中，我们并不知道真实的 \(g(X_t)\)，也不知道哪些变量真正具有预测能力。

因此，一个 Financial Research Agent 的主要工作通常包括两部分：
\[\underbrace{\text{Discover useful predictors}}_{\text{Factor Research}}
+
\underbrace{\text{Combine predictors effectively}}_{\text{Model Research}}.\]

本阶段，我们将从实际金融研究流程出发，理解 Agent 究竟在发现、修改和优化什么。

\hypertarget{part-4-63-0001-learning-roadmap}{%
\subsection{0. Learning Roadmap}\label{part-4-63-0001-learning-roadmap}}

\begin{readertable}
\begin{longtable}[]{@{}
  >{\raggedright\arraybackslash}p{(\columnwidth - 4\tabcolsep) * \real{0.36}}
  >{\raggedright\arraybackslash}p{(\columnwidth - 4\tabcolsep) * \real{0.44}}
  >{\raggedright\arraybackslash}p{(\columnwidth - 4\tabcolsep) * \real{0.16}}@{}}
\toprule
Topic & Core Question & Priority \\
\midrule
\endhead
Alpha \& Beta & 金融研究中的 Alpha 究竟是什么？ & P0 \\
Financial Factors & 因子与普通 Feature 有什么区别？ & P0 \\
Cross-sectional Prediction & 如何在多个资产之间进行相对预测？ & P0 \\
IC \& Rank IC & 如何评价因子的预测能力？ & P0 \\
Factor Zoo & 大量已发现的因子是否真实有效？ & P0 \\
Regularized Regression & 如何控制预测模型的复杂度？ & P0 \\
Nonlinear ML Models & 非线性模型能够带来什么？ & P1 \\
Automated Factor Mining & Agent 如何自动发现和改进金融因子？ & P0 \\
\bottomrule
\end{longtable}
\end{readertable}

\hypertarget{part-4-63-0001-quantitative-research-pipeline}{%
\subsection{1. Quantitative Research Pipeline}\label{part-4-63-0001-quantitative-research-pipeline}}

首先建立整个金融预测研究的基本架构。

\bookgraphic{0.290}{part-4-63-0001-001-01}{简化的 Quantitative Research Pipeline。蓝色模块是本阶段的重点，Portfolio Construction 和经济绩效评价将在 IV-D 展开。}

对于普通的金融机器学习项目，人类研究人员负责决定如何修改 Factor Engineering 和 Prediction Model。

对于 Automated Financial Research，Agent 可以自动完成这些修改。

对于 RSI，系统还可能进一步修改负责发现因子、选择模型、组织实验的研究机制本身。

因此，金融模型的演化与研究 Agent 的演化是两个不同层次。

\hypertarget{part-4-63-0001-alpha-and-betaux91d1ux878dux7814ux7a76ux4e2dux7684ux4e24ux4e2aux6838ux5fc3ux6982ux5ff5}{%
\subsection{2. Alpha and Beta（金融研究中的两个核心概念）}\label{part-4-63-0001-alpha-and-betaux91d1ux878dux7814ux7a76ux4e2dux7684ux4e24ux4e2aux6838ux5fc3ux6982ux5ff5}}

\hypertarget{part-4-63-0001-market-model}{%
\subsubsection{2.1 Market Model}\label{part-4-63-0001-market-model}}

首先考虑一个经典的线性资产收益模型。

为了区分不同资产，引入下标 \(i\) 表示第 \(i\) 个资产，继续使用 \(t\) 表示金融时间。

定义：

\begin{itemize}
\tightlist
\item
  \(r_{i,t}\)：资产 \(i\) 在时刻 \(t\) 的收益率；
\item
  \(r_{m,t}\)：市场组合收益率；
\item
  \(r_{f,t}\)：无风险收益率。
\end{itemize}

建立如下回归模型：
\[\boxed{
r_{i,t}-r_{f,t}
=
\alpha_i+\beta_i(r_{m,t}-r_{f,t})
+\varepsilon_{i,t}
}\]

其中，\(\alpha_i\) 和 \(\beta_i\) 分别是回归截距和斜率。

\hypertarget{part-4-63-0001-beta-systematic-risk-exposure}{%
\paragraph{Beta --- Systematic Risk Exposure}\label{part-4-63-0001-beta-systematic-risk-exposure}}

\[\boxed{
\beta_i=
\frac{
\operatorname{Cov}(r_{i,t}-r_{f,t},r_{m,t}-r_{f,t})
}{
\operatorname{Var}(r_{m,t}-r_{f,t})
}
}\]

Beta 衡量资产对市场超额收益变化的敏感程度。

\textbf{Understanding Market Beta}

Market exposure \(\beta_i\)\\
\textbf{1.2}

1.2（范围 -1--3）

Market excess return\\
\textbf{5.0\%}

5（范围 -10--10）

\bookgraphic{0.618}{part-4-63-0001-002-01}{}

\[\beta_i(r_{m,t}-r_{f,t})=6.00\%\]

仅计算市场风险暴露所对应的收益成分，不包含 Alpha 与随机误差，也不是实际收益预测。

beta = -1

\bookgraphic{0.618}{part-4-63-0001-002-state-01}{}

通常：

\begin{readertable}
\begin{longtable}[]{@{}
  >{\raggedleft\arraybackslash}p{(\columnwidth - 2\tabcolsep) * \real{0.21}}
  >{\raggedright\arraybackslash}p{(\columnwidth - 2\tabcolsep) * \real{0.75}}@{}}
\toprule
Beta & Interpretation \\
\midrule
\endhead
0 & 不存在对该市场因子的线性暴露 \\
0.5 & 市场敏感程度较低 \\
1 & 市场敏感程度与基准相当 \\
2 & 市场敏感程度较高 \\
-1 & 与市场变化呈反向线性关系 \\
\bottomrule
\end{longtable}
\end{readertable}

Beta 是线性风险暴露指标，不等于资产的全部风险。

\hypertarget{part-4-63-0001-alpha-risk-adjusted-excess-performance}{%
\paragraph{Alpha --- Risk-adjusted Excess Performance}\label{part-4-63-0001-alpha-risk-adjusted-excess-performance}}

如果回归残差均值为零，则：
\[\boxed{
\alpha_i
=
\mathbb E[r_{i,t}-r_{f,t}]
-
\beta_i\mathbb E[r_{m,t}-r_{f,t}]
}\]

Alpha 衡量不能被指定风险模型解释的平均超额收益。

例如，假设某一期：

\begin{readertable}
\begin{longtable}[]{@{}
  >{\raggedright\arraybackslash}p{(\columnwidth - 2\tabcolsep) * \real{0.74}}
  >{\raggedleft\arraybackslash}p{(\columnwidth - 2\tabcolsep) * \real{0.22}}@{}}
\toprule
Item & Value \\
\midrule
\endhead
Asset Return & 12\% \\
Market Return & 10\% \\
Risk-free Return & 2\% \\
Beta & 1.2 \\
Market-model Benchmark Return & 11.6\% \\
Realized Abnormal Return & 0.4\% \\
\bottomrule
\end{longtable}
\end{readertable}

虽然该资产比市场多获得了 2 个百分点的收益，但经过市场风险暴露调整之后，该期异常收益只有 0.4 个百分点。

单期异常收益并不等于统计意义上的 Alpha，后者通常需要通过多个时期的数据进行估计。

\hypertarget{part-4-63-0001-multi-factor-modelux591aux56e0ux5b50ux6a21ux578b}{%
\subsubsection{2.2 Multi-factor Model（多因子模型）}\label{part-4-63-0001-multi-factor-modelux591aux56e0ux5b50ux6a21ux578b}}

市场收益率并不是唯一可能解释资产收益的变量。

因此，可以将 Market Model 推广为：
\[\boxed{
r_{i,t}-r_{f,t}
=
\alpha_i+
\sum_{j=1}^{p}\beta_{i,j}F_{j,t}
+\varepsilon_{i,t}
}\]

其中：

\(F_{j,t}\) 表示第 \(j\) 个风险因子的实现值，\(\beta_{i,j}\) 表示资产 \(i\) 对该因子的风险暴露。

经典的 Fama--French Three-factor Model 使用三个股票市场因子：
\[\begin{aligned}
r_{i,t}-r_{f,t}
={}&\alpha_i
+\beta_{i,m}(r_{m,t}-r_{f,t})\\
&+\beta_{i,s}\mathrm{SMB}_t\\
&+\beta_{i,v}\mathrm{HML}_t
+\varepsilon_{i,t}.
\end{aligned}\]

\begin{readertable}
\begin{longtable}[]{@{}
  >{\raggedright\arraybackslash}p{(\columnwidth - 4\tabcolsep) * \real{0.14}}
  >{\raggedright\arraybackslash}p{(\columnwidth - 4\tabcolsep) * \real{0.31}}
  >{\raggedright\arraybackslash}p{(\columnwidth - 4\tabcolsep) * \real{0.51}}@{}}
\toprule
Factor & Full Name & Economic Meaning \\
\midrule
\endhead
MKT & Market Excess Return & 市场风险暴露 \\
SMB & Small Minus Big & 小市值与大市值股票组合之间的收益差 \\
HML & High Minus Low & 高与低账面市值比股票组合之间的收益差 \\
\bottomrule
\end{longtable}
\end{readertable}

这三个因子来源于 Fama 和 French 1993 年发表的经典资产定价研究。

\href{https://doi.org/10.1016/0304-405X\%2893\%2990023-5}{来源 1}

这里还有一个值得关注的结论：

\textbf{Alpha 是相对于某个 Risk Model 定义的，并不是一个独立于模型的绝对量。}

在单因子模型中观察到的 Alpha，加入其他风险因子之后可能显著缩小。

这对 Financial RSI 非常重要。Agent 可能发现了一个能够提高历史收益的方法，但其收益增加实际上来自更大的市场风险暴露，而不是真正的预测能力改善。

\hypertarget{part-4-63-0001-financial-factorsux91d1ux878dux56e0ux5b50}{%
\subsection{3. Financial Factors（金融因子）}\label{part-4-63-0001-financial-factorsux91d1ux878dux56e0ux5b50}}

\hypertarget{part-4-63-0001-factor-ux4e0e-feature-ux7684ux533aux522b}{%
\subsubsection{3.1 Factor 与 Feature 的区别}\label{part-4-63-0001-factor-ux4e0e-feature-ux7684ux533aux522b}}

这是金融机器学习中一个比较容易混淆的术语问题。

Factor 在不同研究领域中至少有两种常见含义。

\begin{readertable}
\begin{longtable}[]{@{}
  >{\raggedright\arraybackslash}p{(\columnwidth - 4\tabcolsep) * \real{0.35}}
  >{\raggedright\arraybackslash}p{(\columnwidth - 4\tabcolsep) * \real{0.39}}
  >{\raggedright\arraybackslash}p{(\columnwidth - 4\tabcolsep) * \real{0.22}}@{}}
\toprule
Concept & Meaning & Example \\
\midrule
\endhead
Risk Factor & 解释资产共同风险或风险溢价的变量 & Fama--French SMB \\
Predictive Factor / Alpha Factor & 用于预测未来收益率的资产特征 & 过去 20 日收益率 \\
ML Feature & 输入机器学习模型的变量 & 标准化交易量 \\
Factor Exposure & 资产对某个风险因子的敏感程度 & 市场 Beta \\
\bottomrule
\end{longtable}
\end{readertable}

在现代量化投研系统中，Predictive Factor 和 ML Feature 经常具有相近含义，但它们不能与 Asset Pricing 中的 Risk Factor 完全混为一谈。

例如：

\begin{itemize}
\tightlist
\item
  某只股票的市值是一个 Characteristic。
\item
  根据市值构造的小盘股减大盘股组合收益是一个 Factor Return。
\item
  股票对该 Factor Return 的回归系数是 Factor Exposure。
\end{itemize}

这是三个不同的数学对象。

接下来我们主要研究 Predictive Factor，因为它与你当前准备的 RSI 课题联系更加直接。

\hypertarget{part-4-63-0001-mathematical-representation}{%
\subsubsection{3.2 Mathematical Representation}\label{part-4-63-0001-mathematical-representation}}

沿用 IV-B 中的 DGP，并增加资产索引：
\[r_{i,t+1}
=
g(X_{i,t})+\varepsilon_{i,t+1}.\]

其中：
\[X_{i,t}
=
\begin{pmatrix}
X_{i,t}^{(1)}\\
X_{i,t}^{(2)}\\
\vdots\\
X_{i,t}^{(p)}
\end{pmatrix}
\in\mathbb R^p.\]

每一个 \(X_{i,t}^{(j)}\) 都可以表示一个金融预测因子。

金融预测模型的目标是利用这些因子估计未知函数 \(g\)。

\[\boxed{
\hat r_{i,t+1}
=
h(X_{i,t})
}\]

于是，Factor Research 负责决定 \(X_{i,t}\) 包含什么信息，Model Research 则负责决定 \(h\) 如何利用这些信息。

\hypertarget{part-4-63-0001-common-factor-categories}{%
\subsubsection{3.3 Common Factor Categories}\label{part-4-63-0001-common-factor-categories}}

\bookgraphic{0.391}{part-4-63-0001-004-01}{}

\textbf{Momentum（动量）}

\begin{verbatim}
  过去表现较好的资产，是否具有继续相对表现较好的倾向？
\end{verbatim}

\[X_{i,t}^{\mathrm{mom}}=\frac{P_{i,t}}{P_{i,t-L}}-1.\]

\bookgraphic{0.391}{part-4-63-0001-004-02}{}

\textbf{Reversal（反转）}

\begin{verbatim}
  近期价格变化是否存在短期反向调整？
\end{verbatim}

\[X_{i,t}^{\mathrm{rev}}
      =-r_{i,t}.\]

\bookgraphic{0.391}{part-4-63-0001-004-03}{}

\textbf{Value（价值）}

\begin{verbatim}
  资产价格相对于企业基本面是否较低？
\end{verbatim}

\[X_{i,t}^{\mathrm{value}}
      =
      \frac{\text{Book Equity}_{i,t}}
      {\text{Market Equity}_{i,t}}.\]

\bookgraphic{0.391}{part-4-63-0001-004-04}{}

\textbf{Volatility（波动率）}

\begin{verbatim}
  利用近期收益率的标准差刻画资产风险。
\end{verbatim}

\[X_{i,t}^{\mathrm{vol}}
      =
      \operatorname{Std}(r_{i,t-L+1:t}).\]

\bookgraphic{0.391}{part-4-63-0001-004-05}{}

\textbf{Liquidity \& Volume（流动性与交易量）}

\begin{verbatim}
  交易活跃程度、市场流动性与未来收益是否存在统计关系？
\end{verbatim}

\[X_{i,t}^{\mathrm{volume}}
      =
      \frac{V_{i,t}}{\operatorname{MA}_L(V_{i,t})}.\]

其中 \(L\) 表示 Lookback Window Length，\(V_{i,t}\) 表示交易量。

以上只是因子的基本构造示例，并不意味着它们在任何市场、时期或预测期限下都具有稳定的预测价值。

\hypertarget{part-4-63-0001-cross-sectional-predictionux6a2aux622aux9762ux9884ux6d4b}{%
\subsection{4. Cross-sectional Prediction（横截面预测）}\label{part-4-63-0001-cross-sectional-predictionux6a2aux622aux9762ux9884ux6d4b}}

\hypertarget{part-4-63-0001-time-series-vs.-cross-section}{%
\subsubsection{4.1 Time Series vs. Cross-section}\label{part-4-63-0001-time-series-vs.-cross-section}}

此前 IV-A 主要研究单个资产随时间变化的收益率序列。

但实际金融数据经常同时包含多个资产和多个时间点，从而形成 Panel Data（面板数据）。

\bookgraphic{0.499}{part-4-63-0001-005-01}{横向研究同一资产随时间的变化；纵向研究同一时刻不同资产之间的差异。}

\begin{readertable}
\begin{longtable}[]{@{}
  >{\raggedright\arraybackslash}p{(\columnwidth - 4\tabcolsep) * \real{0.22}}
  >{\raggedright\arraybackslash}p{(\columnwidth - 4\tabcolsep) * \real{0.36}}
  >{\raggedright\arraybackslash}p{(\columnwidth - 4\tabcolsep) * \real{0.39}}@{}}
\toprule
& Time-series Prediction & Cross-sectional Prediction \\
\midrule
\endhead
研究对象 & 同一个资产 & 同一时刻多个资产 \\
核心问题 & 明天市场会上涨还是下跌？ & 明天哪些资产相对表现更好？ \\
主要依赖结构 & Temporal Dependence & Cross-asset Dependence \\
常见指标 & MSE、MAE、OOS \(R^2\) & IC、Rank IC、Portfolio Spread \\
\bottomrule
\end{longtable}
\end{readertable}

这两种预测任务并不互斥。

一个机器学习模型可以利用跨资产、跨时间的 Panel Data 进行训练，然后每天输出横截面预测结果。

\hypertarget{part-4-63-0001-ux4e3aux4ec0ux4e48ux76f8ux5bf9ux9884ux6d4bux4e0eux7eddux5bf9ux9884ux6d4bux4e0dux540c}{%
\subsubsection{4.2 为什么相对预测与绝对预测不同？}\label{part-4-63-0001-ux4e3aux4ec0ux4e48ux76f8ux5bf9ux9884ux6d4bux4e0eux7eddux5bf9ux9884ux6d4bux4e0dux540c}}

假设有两个资产：

\begin{readertable}
\begin{longtable}[]{@{}
  >{\raggedright\arraybackslash}p{(\columnwidth - 4\tabcolsep) * \real{0.19}}
  >{\raggedleft\arraybackslash}p{(\columnwidth - 4\tabcolsep) * \real{0.42}}
  >{\raggedleft\arraybackslash}p{(\columnwidth - 4\tabcolsep) * \real{0.35}}@{}}
\toprule
Asset & Predicted Return & Actual Return \\
\midrule
\endhead
A & 3\% & -2\% \\
B & 1\% & -5\% \\
\bottomrule
\end{longtable}
\end{readertable}

从绝对收益率预测来看，模型对两个资产的预测都存在明显误差。

但从相对排序来看：
\[\hat r_A>\hat r_B,\]

并且：
\[r_A>r_B.\]

因此模型正确识别了两个资产之间的相对表现。

这说明：
\[\boxed{
\text{Accurate Return Ranking}
\neq
\text{Accurate Return Level Prediction}
}\]

所以金融研究不能仅仅依赖 MSE，还需要能够衡量排序质量的统计指标。

\hypertarget{part-4-63-0001-cross-sectional-dependence}{%
\subsubsection{4.3 Cross-sectional Dependence}\label{part-4-63-0001-cross-sectional-dependence}}

另一个重要问题是，Panel Data 中的样本并不必然相互独立。

假设：
\[r_{i,t}=\beta_iF_t+\varepsilon_{i,t},\]

其中 \(F_t\) 是共同市场冲击，残差之间不相关，且残差与 \(F_t\) 不相关。

对于不同资产 \(i\neq j\)：
\[\boxed{
\operatorname{Cov}(r_{i,t},r_{j,t})
=
\beta_i\beta_j\operatorname{Var}(F_t).
}\]

即使各资产的个体误差相互独立，共同市场因子仍然可以导致资产收益率相关。

因此，拥有 500 只股票、1,000 个交易日，并不意味着获得了 500,000 个 IID 样本。

这是 IV-A 中 Effective Sample Size 问题在横截面维度上的进一步推广。

\hypertarget{part-4-63-0001-information-coefficient-ic}{%
\subsection{5. Information Coefficient (IC)}\label{part-4-63-0001-information-coefficient-ic}}

这是本阶段最重要的实用评价指标之一。

\hypertarget{part-4-63-0001-pearson-ic}{%
\subsubsection{5.1 Pearson IC}\label{part-4-63-0001-pearson-ic}}

假设在时刻 \(t\)，我们对 \(N_t\) 只股票计算了某个因子值或者预测得分。

定义：
\[\boxed{
\mathrm{IC}_t
=
\operatorname{Corr}_{i}
\left(
X_{i,t}^{(j)},r_{i,t+1}
\right)
}\]

这里相关系数是沿资产维度 \(i\) 计算，而不是沿时间维度计算。

如果我们希望评价整个预测模型，则将因子值替换为预测值：
\[\boxed{
\mathrm{IC}_t
=
\operatorname{Corr}_{i}
\left(
\hat r_{i,t+1},r_{i,t+1}
\right)
}\]

它实际上就是每一个时刻的 Cross-sectional Pearson Correlation。

\hypertarget{part-4-63-0001-rank-ic}{%
\subsubsection{5.2 Rank IC}\label{part-4-63-0001-rank-ic}}

Rank IC 则使用 Spearman Rank Correlation：
\[\boxed{
\mathrm{RankIC}_t
=
\operatorname{Corr}_{i}
\left(
\operatorname{rank}(\hat r_{i,t+1}),
\operatorname{rank}(r_{i,t+1})
\right)
}\]

在没有并列名次的情况下，可以写成：
\[\boxed{
\mathrm{RankIC}_t
=
1-
\frac{
6\sum_{i=1}^{N_t}d_i^2
}{
N_t(N_t^2-1)
}
}\]

其中 \(d_i\) 表示预测名次与实际名次之间的差。

Microsoft Qlib 等量化研究框架也使用 Pearson IC 和 Spearman Rank IC 作为模型评价指标。

\href{https://qlib.readthedocs.io/en/stable/component/report.html}{来源 1} \href{https://github.com/microsoft/qlib/blob/main/docs/component/report.rst}{来源 2}

\textbf{Rank IC --- Interactive Example}

比较五只股票的预测名次与实际名次，名次 1 表示收益最高。

Perfect ranking

Minor ranking errors

Weak ranking agreement

Reversed ranking

\bookgraphic{0.459}{part-4-63-0001-007-01}{}

Sum of squared rank errors

\textbf{0}

Rank IC

\textbf{1.00}

教学数据，用于理解 Spearman Correlation 对排序错误的反应。

\hypertarget{part-4-63-0001-why-rank-ic}{%
\subsubsection{5.3 Why Rank IC?}\label{part-4-63-0001-why-rank-ic}}

Rank IC 具有一个非常重要的数学性质：对严格单调递增变换保持不变。

假设：
\[z_i=f(x_i),\]

其中 \(f\) 严格单调递增。

那么：
\[\operatorname{rank}(z_i)
=
\operatorname{rank}(x_i).\]

因此：
\[\boxed{
\rho_{\mathrm{Spearman}}(f(X),Y)
=
\rho_{\mathrm{Spearman}}(X,Y)
}\]

但 Pearson Correlation 对一般非线性单调变换并不保持不变。

\textbf{Linear Relationship}

\bookgraphic{0.341}{part-4-63-0001-008-01}{Pearson IC = 1}

Rank IC = 1

\textbf{Nonlinear Monotonic Relationship}

\bookgraphic{0.341}{part-4-63-0001-008-02}{Pearson IC 小于 1}

Rank IC = 1

Rank IC 更注重相对排序，而不是预测数值的线性校准。

但它也会丢失 Magnitude Information：预测收益为 1\% 还是 100\%，只要排序不变，Rank IC 就可能完全相同。

所以，Rank IC 高并不自动意味着预测值准确，更不意味着交易后的经济收益一定较高。

\hypertarget{part-4-63-0001-mean-ic-and-icir}{%
\subsubsection{5.4 Mean IC and ICIR}\label{part-4-63-0001-mean-ic-and-icir}}

对于 \(T\) 个评价时期，可以定义平均 IC：
\[\boxed{
\overline{\mathrm{IC}}
=
\frac1T\sum_{t=1}^{T}\mathrm{IC}_t
}\]

进一步定义 Information Coefficient Information Ratio：
\[\boxed{
\mathrm{ICIR}
=
\frac{
\overline{\mathrm{IC}}
}{
s_{\mathrm{IC}}
}
}\]

其中 \(s_{\mathrm{IC}}\) 为 IC 序列的样本标准差，假设其大于零。

\bookgraphic{0.786}{part-4-63-0001-009-01}{}

假想因子表现。Factor A 偶尔具有较高 IC，但波动较大；Factor B 的 IC 较低，却更加稳定。

\begin{readertable}
\begin{longtable}[]{@{}
  >{\raggedright\arraybackslash}p{(\columnwidth - 2\tabcolsep) * \real{0.40}}
  >{\raggedright\arraybackslash}p{(\columnwidth - 2\tabcolsep) * \real{0.56}}@{}}
\toprule
Metric & What It Measures \\
\midrule
\endhead
Mean IC & 平均预测相关性 \\
IC Standard Deviation & 预测相关性的波动 \\
ICIR & 平均预测相关性相对于波动的大小 \\
IC Positive Rate & IC 为正的时期所占比例 \\
\bottomrule
\end{longtable}
\end{readertable}

ICIR 本身不是 \(t\)-statistic。

只有在适当的 IID 假设下，才有：
\[t_{\mathrm{IC}}
=
\sqrt T\cdot\mathrm{ICIR}.\]

如果 IC 序列存在自相关、重叠标签或结构变化，则需要考虑 HAC Standard Error 或 Block Bootstrap。

这是 IV-A 中时间依赖问题的又一次实际应用。

\hypertarget{part-4-63-0001-the-mathematical-relationship-between-ic-and-r2}{%
\subsection{\texorpdfstring{6. The Mathematical Relationship between IC and \(R^2\)}{6. The Mathematical Relationship between IC and R\^{}2}}\label{part-4-63-0001-the-mathematical-relationship-between-ic-and-r2}}

这一节很值得深入研究，因为它能够解释 IC 和 Regression Performance 之间的本质关系。

\hypertarget{part-4-63-0001-problem-setting}{%
\subsubsection{6.1 Problem Setting}\label{part-4-63-0001-problem-setting}}

为简化公式，本节暂时使用：
\[Y=r_{t+1},
\qquad
Z=h(X_t).\]

其中 \(Z\) 是一个已经确定的预测得分。

假设两者均具有有限且非零的方差。

我们考虑对预测得分进行 Linear Calibration：
\[\tilde h(Z)=a+bZ.\]

目标是找到：
\[(a^*,b^*)
=
\arg\min_{a,b}
\mathbb E[(Y-a-bZ)^2].\]

\hypertarget{part-4-63-0001-theorem-optimal-linear-calibration}{%
\subsubsection{6.2 Theorem --- Optimal Linear Calibration}\label{part-4-63-0001-theorem-optimal-linear-calibration}}

根据 Population OLS 的结论：
\[b^*
=
\frac{\operatorname{Cov}(Y,Z)}
{\operatorname{Var}(Z)},\]

\[a^*
=
\mathbb E[Y]-b^*\mathbb E[Z].\]

代入最优平方预测风险：
\[R(\tilde h^*)
=
\operatorname{Var}(Y)
-
\frac{\operatorname{Cov}(Y,Z)^2}
{\operatorname{Var}(Z)}.\]

于是：
\[\begin{aligned}
R^2(\tilde h^*)
&=
1-\frac{R(\tilde h^*)}{\operatorname{Var}(Y)}\\
&=
\frac{
\operatorname{Cov}(Y,Z)^2
}{
\operatorname{Var}(Y)\operatorname{Var}(Z)
}\\
&=
\operatorname{Corr}(Y,Z)^2.
\end{aligned}\]

最终：
\[\boxed{
R^2(\tilde h^*)=\rho_{Y,Z}^2
}\]

\(\square\)

这意味着，Population Pearson Correlation 的平方等于对该预测得分进行最优仿射校准之后所获得的 Population \(R^2\)。

但请注意三个限定条件：需要最优线性校准、需要采用总体均值基准、讨论的是 Pearson Correlation 而非 Rank Correlation。

不能直接宣称：
\[R^2(h)=\mathrm{IC}^2.\]

特别是跨日期平均 IC 的平方，不等于整个金融面板数据的 OOS \(R^2\)。

\hypertarget{part-4-63-0001-maximum-achievable-correlation}{%
\subsubsection{6.3 Maximum Achievable Correlation}\label{part-4-63-0001-maximum-achievable-correlation}}

现在结合 IV-B 中的 Bayes Optimal Predictor。

对于任意有限方差、非常数预测函数 \(h(X_t)\)，由条件期望的正交性质：
\[\operatorname{Cov}(r_{t+1},h(X_t))
=
\operatorname{Cov}(g(X_t),h(X_t)).\]

利用 Cauchy--Schwarz Inequality：
\[\operatorname{Cov}^2(g(X_t),h(X_t))
\leq
\operatorname{Var}(g(X_t))
\operatorname{Var}(h(X_t)).\]

所以：
\[\begin{aligned}
\operatorname{Corr}^2(r_{t+1},h(X_t))
&\leq
\frac{
\operatorname{Var}(g(X_t))
}{
\operatorname{Var}(r_{t+1})
}\\
&=
\frac{\mathrm{SNR}}{1+\mathrm{SNR}}.
\end{aligned}\]

因此：
\[\boxed{
|\operatorname{Corr}(r_{t+1},h(X_t))|
\leq
\sqrt{\frac{\mathrm{SNR}}{1+\mathrm{SNR}}}
}\]

当 \(g(X_t)\) 具有非零方差时，取 \(h=g\) 可以达到这一上界。

\textbf{SNR --- Theoretical Prediction Limits}

Signal-to-noise ratio\\
\textbf{0.0100}

-2（范围 -4--0）

Maximum Population R²

\textbf{0.990\%}

Maximum Absolute Correlation

\textbf{0.0995}

基于 IV-B 的方差信噪比定义。这里讨论总体 Pearson 相关系数，而非有限样本的平均横截面 IC。

这个结果也解释了为什么在低 SNR 环境中，即使较小的预测相关性也可能具有研究价值。

例如：

\begin{readertable}
\begin{longtable}[]{@{}
  >{\raggedleft\arraybackslash}p{(\columnwidth - 4\tabcolsep) * \real{0.24}}
  >{\raggedleft\arraybackslash}p{(\columnwidth - 4\tabcolsep) * \real{0.33}}
  >{\raggedleft\arraybackslash}p{(\columnwidth - 4\tabcolsep) * \real{0.39}}@{}}
\toprule
Task-level SNR & Optimal Population \(R^2\) & Maximum Absolute Correlation \\
\midrule
\endhead
0.001 & 0.10\% & 0.0316 \\
0.01 & 0.99\% & 0.0995 \\
0.1 & 9.09\% & 0.3015 \\
1 & 50\% & 0.7071 \\
\bottomrule
\end{longtable}
\end{readertable}

\hypertarget{part-4-63-0001-the-factor-zoo-problem}{%
\subsection{7. The Factor Zoo Problem}\label{part-4-63-0001-the-factor-zoo-problem}}

这是金融因子研究与 RSI 最直接的理论交汇点。

\hypertarget{part-4-63-0001-what-is-factor-zoo}{%
\subsubsection{7.1 What Is Factor Zoo?}\label{part-4-63-0001-what-is-factor-zoo}}

数十年来，金融研究人员不断提出新的因子，试图解释资产收益率的横截面差异。

这些因子的构造方式包括价格趋势、估值指标、企业规模、盈利能力、交易量、波动率以及更复杂的变量组合。

由此产生了所谓的 \textbf{Factor Zoo（因子动物园）}：大量被报告为有效的金融因子，其中可能存在冗余、虚假发现和不稳定的预测关系。

Harvey、Liu 和 Zhu 在 2016 年发表于 \emph{The Review of Financial Studies} 的研究中指出，在长期、大规模的因子检验背景下，传统的显著性标准可能过于宽松。他们的模型建议，新因子需要达到高于 3 的 \(t\)-statistic 门槛。

\href{https://academic.oup.com/rfs/article-abstract/29/1/5/1843824?login=false}{来源 1} \href{https://www.nber.org/papers/w20592}{来源 2}

这不是一个适用于所有金融检验的统一阈值，而是特定多重检验分析下得到的研究结论。

\hypertarget{part-4-63-0001-rsi-ux4e3aux4ec0ux4e48ux53efux80fdux52a0ux5267-factor-zoo}{%
\subsubsection{7.2 RSI 为什么可能加剧 Factor Zoo？}\label{part-4-63-0001-rsi-ux4e3aux4ec0ux4e48ux53efux80fdux52a0ux5267-factor-zoo}}

传统人类研究与自动化 Agent Research 的一个重要差别在于实验规模。

\textbf{Hypothetical Factor Discovery Scale}

Manual research

\begin{verbatim}
    **20**
\end{verbatim}

Automated factor generation

\begin{verbatim}
    **200**
\end{verbatim}

Large-scale adaptive search

\begin{verbatim}
    **2,000**
\end{verbatim}

Repeated recursive optimization

\begin{verbatim}
    **20,000**
\end{verbatim}

假想候选因子数量，条形长度采用对数尺度，仅说明自动化搜索可能扩大研究空间，不代表实际系统的实验规模。

结合 II-C：

如果有 \(m\) 个完全无效的因子，每个因子对应一个有效的显著性水平为 \(\alpha\) 的检验，那么错误拒绝次数 \(V\) 的期望为：
\[\mathbb E[V]=m\alpha\]

其中假设全部零假设为真且各检验达到名义错误率，不需要检验之间相互独立。

\textbf{Factor Zoo --- False Discovery Explorer}

Number of tested noise factors 100 / 1,000 / 10,000 Individual significance level

5\%

1\%

0.1\%

Expected False Positives

\textbf{50}

Bonferroni threshold for 5\% FWER

\textbf{0.000050}

假设所有候选因子均无效且单项检验有效。真实 RSI 的自适应假设生成可能使普通 p-value 本身失效。

这就是为什么自动化因子生成系统必须同时研究 Validation Mechanism。

否则，研究效率提高可能意味着发现虚假因子的效率也提高了。

\hypertarget{part-4-63-0001-factor-redundancy}{%
\subsubsection{7.3 Factor Redundancy}\label{part-4-63-0001-factor-redundancy}}

Factor Zoo 的另一个重要问题并不是因子完全无效，而是新增因子没有提供独立于已有因子的 Incremental Information。

假设已有因子：
\[X^{(1)},\ldots,X^{(p)}.\]

新增因子为：
\[X^{(p+1)}
=
f(X^{(1)},\ldots,X^{(p)}).\]

根据 IV-B 中 Information Set 的结论，这种确定性变换不会增加总体信息集合。

但它可能使有限模型更容易利用非线性结构。

另一方面，大量近似重复的因子也可能导致 Multicollinearity 和估计不稳定。

Feng、Giglio 和 Xiu 在 2020 年发表的 \emph{Taming the Factor Zoo: A Test of New Factors} 就研究了如何在大量既有因子的背景下评价新增因子的边际贡献。

\href{https://onlinelibrary.wiley.com/doi/full/10.1111/jofi.12883}{来源 1}

\hypertarget{part-4-63-0001-regularized-regressionux6b63ux5219ux5316ux56deux5f52}{%
\subsection{8. Regularized Regression（正则化回归）}\label{part-4-63-0001-regularized-regressionux6b63ux5219ux5316ux56deux5f52}}

这是本阶段第二个需要重点掌握的数学部分。

\hypertarget{part-4-63-0001-ordinary-least-squares}{%
\subsubsection{8.1 Ordinary Least Squares}\label{part-4-63-0001-ordinary-least-squares}}

定义训练数据：
\[X\in\mathbb R^{n\times p},
\qquad
\boldsymbol y\in\mathbb R^n.\]

其中 \(n\) 为样本数，\(p\) 为特征数。

假设数据已经中心化，暂时省略截距。

OLS 的目标函数为：
\[\boxed{
\hat{\boldsymbol\beta}_{\mathrm{OLS}}
=
\arg\min_{\boldsymbol\beta}
\|\boldsymbol y-X\boldsymbol\beta\|_2^2
}\]

如果 \(X^\top X\) 可逆，则：
\[\boxed{
\hat{\boldsymbol\beta}_{\mathrm{OLS}}
=
(X^\top X)^{-1}X^\top\boldsymbol y
}\]

但当因子数量很多、因子之间存在高度相关性时，\(X^\top X\) 可能接近奇异矩阵。

此时即使存在逆矩阵，估计结果也可能非常不稳定。

\hypertarget{part-4-63-0001-ridge-regression-l_2-regularization}{%
\subsubsection{\texorpdfstring{8.2 Ridge Regression --- \(L_2\) Regularization}{8.2 Ridge Regression --- L\_2 Regularization}}\label{part-4-63-0001-ridge-regression-l_2-regularization}}

Ridge Regression 在 OLS 的基础上加入参数范数惩罚：
\[\boxed{
\hat{\boldsymbol\beta}_{\mathrm{Ridge}}
=
\arg\min_{\boldsymbol\beta}
\left[
\|\boldsymbol y-X\boldsymbol\beta\|_2^2
+
\lambda\|\boldsymbol\beta\|_2^2
\right]
}\]

其中 \(\lambda\geq0\)。

\hypertarget{part-4-63-0001-closed-form-solution}{%
\paragraph{Closed-form Solution}\label{part-4-63-0001-closed-form-solution}}

定义：
\[L(\boldsymbol\beta)
=
(\boldsymbol y-X\boldsymbol\beta)^\top
(\boldsymbol y-X\boldsymbol\beta)
+
\lambda\boldsymbol\beta^\top\boldsymbol\beta.\]

对 \(\boldsymbol\beta\) 求梯度：
\[\nabla_{\boldsymbol\beta}L
=
-2X^\top\boldsymbol y
+
2X^\top X\boldsymbol\beta
+
2\lambda\boldsymbol\beta.\]

令梯度为零：
\[(X^\top X+\lambda I_p)\boldsymbol\beta
=
X^\top\boldsymbol y.\]

对于 \(\lambda>0\)，最终得到：
\[\boxed{
\hat{\boldsymbol\beta}_{\mathrm{Ridge}}
=
(X^\top X+\lambda I_p)^{-1}
X^\top\boldsymbol y
}\]

这可能是目前金融机器学习中最值得你优先掌握的模型之一，因为它直接连接了 Linear Algebra、Regression 和 Statistical Generalization。

\hypertarget{part-4-63-0001-why-does-ridge-improve-numerical-stability}{%
\subsubsection{8.3 Why Does Ridge Improve Numerical Stability?}\label{part-4-63-0001-why-does-ridge-improve-numerical-stability}}

考虑 Eigendecomposition：
\[X^\top X=QDQ^\top,\]

其中：
\[D=\operatorname{diag}(d_1,\ldots,d_p),
\qquad d_j\geq0.\]

OLS 涉及：
\[(X^\top X)^{-1}
=
Q\operatorname{diag}
\left(
\frac1{d_1},\ldots,\frac1{d_p}
\right)Q^\top.\]

当某个 \(d_j\) 接近零时，\(1/d_j\) 会非常大。

而 Ridge 将其替换为：
\[(X^\top X+\lambda I_p)^{-1}
=
Q\operatorname{diag}
\left(
\frac1{d_1+\lambda},\ldots,
\frac1{d_p+\lambda}
\right)Q^\top.\]

因此，小特征值方向受到更强的相对抑制。

\textbf{Ridge Regularization --- Eigenvalue Perspective}

Regularization strength \(\lambda\)\\
\textbf{0.50}

0.5（范围 0--5）

\begin{readertable}
\begin{longtable}[]{@{}
  >{\raggedright\arraybackslash}p{(\columnwidth - 4\tabcolsep) * \real{0.31}}
  >{\raggedright\arraybackslash}p{(\columnwidth - 4\tabcolsep) * \real{0.31}}
  >{\raggedright\arraybackslash}p{(\columnwidth - 4\tabcolsep) * \real{0.35}}@{}}
\toprule
Eigenvalue & OLS inverse & Ridge inverse \\
\midrule
\endhead
10.00 & 0.10 & 0.10 \\
2.00 & 0.50 & 0.40 \\
0.50 & 2.00 & 1.00 \\
0.10 & 10.00 & 1.67 \\
0.01 & 100.00 & 1.96 \\
\bottomrule
\end{longtable}
\end{readertable}

Original condition number

\textbf{1,000}

Regularized condition number

\textbf{20.59}

假想特征值示例。条件数对应 \(X^\top X+\lambda I_p\)，并非模型预测误差。

在 Homoskedastic Linear Model 的条件下，Ridge 还会改变参数估计的 Bias--Variance Trade-off。

\[\operatorname{Bias}(\hat{\boldsymbol\beta}_{\mathrm{Ridge}}\mid X)
=
-\lambda(X^\top X+\lambda I_p)^{-1}
\boldsymbol\beta.\]

同时，其条件协方差矩阵为：
\[\begin{aligned}
&\operatorname{Cov}(
\hat{\boldsymbol\beta}_{\mathrm{Ridge}}\mid X)\\
&=
\sigma_\varepsilon^2
(X^\top X+\lambda I_p)^{-1}
X^\top X
(X^\top X+\lambda I_p)^{-1}.
\end{aligned}\]

Ridge 通过引入一定估计偏差来降低估计方差，但是否能改善整体 OOS Risk 取决于真实参数、数据分布和 \(\lambda\) 的选择。

\hypertarget{part-4-63-0001-lasso-l_1-regularization}{%
\subsubsection{\texorpdfstring{8.4 Lasso --- \(L_1\) Regularization}{8.4 Lasso --- L\_1 Regularization}}\label{part-4-63-0001-lasso-l_1-regularization}}

Lasso 将平方范数替换为绝对值范数：
\[\boxed{
\hat{\boldsymbol\beta}_{\mathrm{Lasso}}
=
\arg\min_{\boldsymbol\beta}
\left[
\|\boldsymbol y-X\boldsymbol\beta\|_2^2
+
\lambda\|\boldsymbol\beta\|_1
\right]
}\]

其中：
\[\|\boldsymbol\beta\|_1
=
\sum_{j=1}^{p}|\beta_j|.\]

Lasso 的重要特点是能够将部分系数直接压缩至零，从而实现 Variable Selection。

\textbf{Ridge --- \(L_2\)}

\bookgraphic{0.301}{part-4-63-0001-017-01}{圆形约束边界，通常得到非零但缩小的系数。}

\textbf{Lasso --- \(L_1\)}

\bookgraphic{0.301}{part-4-63-0001-017-02}{菱形约束边界的尖角有利于产生稀疏解。}

正则化的经典几何直觉示意。圆形和菱形分别代表两种范数约束，椭圆代表经验损失等高线；并非具体数值优化结果。

\hypertarget{part-4-63-0001-a-nontrivial-mathematical-result-soft-thresholding}{%
\paragraph{A Nontrivial Mathematical Result --- Soft Thresholding}\label{part-4-63-0001-a-nontrivial-mathematical-result-soft-thresholding}}

考虑一种特别简单的情况：
\[X^\top X=I_p.\]

令：
\[\boldsymbol z=X^\top\boldsymbol y.\]

对于每一个参数坐标，Ridge 的解为：
\[\boxed{
\hat\beta_j^{\mathrm{Ridge}}
=
\frac{z_j}{1+\lambda}
}\]

而 Lasso 的解为：
\[\boxed{
\hat\beta_j^{\mathrm{Lasso}}
=
\operatorname{sign}(z_j)
\left(|z_j|-\frac{\lambda}{2}\right)_+
}\]

这里的 \(\lambda/2\) 来自本节使用的目标函数尺度；如果平方损失前加入 \(1/2\)，阈值也会随之改变。

\textbf{Ridge vs. Lasso --- Coefficient Shrinkage}

Regularization strength\\
\textbf{0.80}

0.8（范围 0--3）

\bookgraphic{0.444}{part-4-63-0001-019-01}{}

OLS

Ridge

Lasso

正交设计矩阵下的精确解。Lasso 将阈值范围内的弱系数直接压缩为零。

l = 3

\bookgraphic{0.444}{part-4-63-0001-019-state-01}{}

这非常适合解释 Lasso 在 Factor Selection 中的作用。

低 SNR 环境中，很多观测系数可能来自噪声，Lasso 提供了一种限制弱信号进入最终模型的方法。

不过，真实信号也可能很弱，因此过强的正则化同样可能删除有价值的因子。

\hypertarget{part-4-63-0001-elastic-net}{%
\subsubsection{8.5 Elastic Net}\label{part-4-63-0001-elastic-net}}

Elastic Net 将两种正则化结合：
\[\boxed{
\begin{aligned}
\hat{\boldsymbol\beta}_{\mathrm{EN}}
=\arg\min_{\boldsymbol\beta}
\big[
&\|\boldsymbol y-X\boldsymbol\beta\|_2^2\\
&+\lambda_1\|\boldsymbol\beta\|_1\\
&+\lambda_2\|\boldsymbol\beta\|_2^2
\big].
\end{aligned}
}\]

\begin{readertable}
\begin{longtable}[]{@{}
  >{\raggedright\arraybackslash}p{(\columnwidth - 6\tabcolsep) * \real{0.19}}
  >{\raggedright\arraybackslash}p{(\columnwidth - 6\tabcolsep) * \real{0.15}}
  >{\raggedright\arraybackslash}p{(\columnwidth - 6\tabcolsep) * \real{0.27}}
  >{\raggedright\arraybackslash}p{(\columnwidth - 6\tabcolsep) * \real{0.34}}@{}}
\toprule
Method & Penalty & Feature Selection & Typical Advantage \\
\midrule
\endhead
OLS & None & No & 简单、可解释 \\
Ridge & \(L_2\) & 通常不直接删除 & 缓解共线性与估计不稳定 \\
Lasso & \(L_1\) & Yes & 稀疏性 \\
Elastic Net & \(L_1+L_2\) & Yes & 兼顾稀疏与相关特征 \\
\bottomrule
\end{longtable}
\end{readertable}

对于 Financial RSI，这些模型不仅可以作为预测方法，还可以成为研究 Overfitting Mechanism 的可控实验对象。

\hypertarget{part-4-63-0001-nonlinear-machine-learning-models}{%
\subsection{9. Nonlinear Machine Learning Models}\label{part-4-63-0001-nonlinear-machine-learning-models}}

线性模型假设：
\[h(X_{i,t})
=
\beta_0+
\sum_{j=1}^{p}\beta_jX_{i,t}^{(j)}.\]

但真实金融预测关系可能存在非线性和特征交互。

例如：
\[g(X^{(1)},X^{(2)})
=
\beta X^{(1)}X^{(2)}.\]

这意味着一个因子的预测作用可能依赖另一个因子的取值。

如果两个因子相互独立、均值均为零，那么：
\[\operatorname{Cov}(g(X^{(1)},X^{(2)}),X^{(1)})=0,\]

另一个因子同样如此。

但是：
\[\operatorname{Var}(g(X^{(1)},X^{(2)}))
=
\beta^2
\operatorname{Var}(X^{(1)})
\operatorname{Var}(X^{(2)}),\]

在非退化情况下可以大于零。

所以，单独研究每一个变量的线性相关性，可能完全遗漏真实的 Interaction Signal。

\textbf{Nonlinear Interaction Signal}

\bookgraphic{0.382}{part-4-63-0001-020-01}{}

Positive signal

Negative signal

合成函数 \(g(X_1,X_2)=X_1X_2\)。单个特征可以不存在边际线性相关性，但两个特征的组合仍然具有预测价值。

这就是 Nonlinear Models 的理论价值。

\hypertarget{part-4-63-0001-common-models}{%
\subsubsection{9.1 Common Models}\label{part-4-63-0001-common-models}}

\begin{readertable}
\begin{longtable}[]{@{}
  >{\raggedright\arraybackslash}p{(\columnwidth - 4\tabcolsep) * \real{0.26}}
  >{\raggedright\arraybackslash}p{(\columnwidth - 4\tabcolsep) * \real{0.36}}
  >{\raggedright\arraybackslash}p{(\columnwidth - 4\tabcolsep) * \real{0.34}}@{}}
\toprule
Model & Fundamental Mechanism & Main Risk \\
\midrule
\endhead
Decision Tree & 递归划分特征空间 & 过度划分数据 \\
Random Forest & 聚合多个随机化决策树 & 噪声特征、分布变化 \\
Gradient Boosting & 逐步拟合残差或损失梯度 & 过多迭代可能导致过拟合 \\
Neural Network & 通过多层非线性变换拟合函数 & 高估计复杂度与训练不稳定 \\
\bottomrule
\end{longtable}
\end{readertable}

对于当前课题，暂时没有必要深入 Backpropagation 或复杂网络结构。

我们只需要理解：它们扩大了模型可以表达的函数集合，可能降低 Approximation Error，但在有限样本条件下也可能增加 Estimation Error。

\hypertarget{part-4-63-0001-actual-financial-evidence}{%
\subsubsection{9.2 Actual Financial Evidence}\label{part-4-63-0001-actual-financial-evidence}}

Gu、Kelly 和 Xiu 在 2020 年发表的 \emph{Empirical Asset Pricing via Machine Learning} 中，对多种金融机器学习方法进行了比较。

\textbf{Monthly Stock-level OOS R²}

Benchmark linear model

\begin{verbatim}
    **0.16%**
\end{verbatim}

Trees / neural networks --- lower end

\begin{verbatim}
    **0.33%**
\end{verbatim}

Trees / neural networks --- upper end

\begin{verbatim}
    **0.40%**
\end{verbatim}

论文报告的部分月度个股收益预测结果，后两项为相关模型得分区间的端点。 \href{https://doi.org/10.1093/rfs/hhaa009}{来源 1}

作者发现，树模型和神经网络能够利用传统模型遗漏的非线性预测变量交互，从而改善其样本外预测表现。

\href{https://www.nber.org/papers/w25398}{来源 1}

这个结果与 IV-B 并不矛盾。

低 SNR 限制理论预测上限，但如果简单模型尚未充分利用已有信息，更复杂的模型仍然可能带来真实改善。

\hypertarget{part-4-63-0001-automated-factor-miningux81eaux52a8ux5316ux56e0ux5b50ux6316ux6398}{%
\subsection{10. Automated Factor Mining（自动化因子挖掘）}\label{part-4-63-0001-automated-factor-miningux81eaux52a8ux5316ux56e0ux5b50ux6316ux6398}}

现在回到 Agentic AI。

在 III-E 中，我们研究过 R\&D-Agent-Quant。它将因子与预测模型的自动化研发组织在同一个协同优化框架中。该论文已发表于 NeurIPS 2025。

\href{https://proceedings.neurips.cc/paper_files/paper/2025/hash/ac5c2b6e423883cbcacbcccf88491b78-Abstract-Datasets_and_Benchmarks_Track.html}{来源 1}

\hypertarget{part-4-63-0001-factor-as-a-program}{%
\subsubsection{10.1 Factor as a Program}\label{part-4-63-0001-factor-as-a-program}}

对于人类研究人员，一个因子可能表现为一条金融公式。

但对于 Research Agent，一个因子完全可以表示为一段程序。

例如：

\begin{verbatim}
def momentum(close, window=20):
    return close / close.shift(window) - 1
\end{verbatim}

Agent 可以进一步修改：

\begin{verbatim}
def volatility_adjusted_momentum(
    close,
    window=20
):
    returns = close.pct_change()

    momentum = (
        close / close.shift(window) - 1
    )

    volatility = (
        returns.rolling(window).std()
    )

    return momentum / (volatility + 1e-8)
\end{verbatim}

从 RSI 的角度来看，这就是一个可被生成、修改、执行和评价的 Code Object。

这里假设输入是一只资产按时间排序的价格序列。实际多资产数据必须分别按资产计算，并核对信息可用时间与成交时间。

\hypertarget{part-4-63-0001-expression-tree-representation}{%
\subsubsection{10.2 Expression Tree Representation}\label{part-4-63-0001-expression-tree-representation}}

因子还可以表示为一个 Expression Tree。

例如：
\[f_t=
\frac{
P_t/P_{t-20}-1
}{
\operatorname{Std}(r_{t-19:t})+\epsilon
}.\]

\bookgraphic{0.375}{part-4-63-0001-024-01}{因子的树结构表示。Agent 可以通过修改运算符、子树、输入变量或窗口长度产生新的候选因子。}

这与你目前学习的 Data Structures \& Algorithm Design 存在非常直接的联系。

Factor Mining 可以被视为在 Expression Space 中进行搜索，常用操作包括 Mutation、Subtree Replacement、Operator Replacement 和 Parameter Adjustment。

\hypertarget{part-4-63-0001-search-space-explosion}{%
\subsubsection{10.3 Search-space Explosion}\label{part-4-63-0001-search-space-explosion}}

假设我们拥有三个基础金融变量：
\[P,\quad V,\quad r.\]

同时拥有一组允许的运算符：
\[+,-,\times,\div,
\operatorname{Mean},
\operatorname{Std},
\operatorname{Rank}.\]

那么，即使原始金融数据的种类有限，也可以生成大量复合因子。

\begin{readertable}
\begin{longtable}[]{@{}
  >{\raggedright\arraybackslash}p{(\columnwidth - 2\tabcolsep) * \real{0.46}}
  >{\raggedright\arraybackslash}p{(\columnwidth - 2\tabcolsep) * \real{0.50}}@{}}
\toprule
Revision & Example \\
\midrule
\endhead
Parameter Mutation & 20 日窗口改为 60 日 \\
Operator Mutation & Moving Mean 改为 Moving Median \\
Feature Combination & Momentum × Volume \\
Nonlinear Transformation & 对因子进行平方 \\
Normalization & Cross-sectional Z-score \\
Risk Adjustment & 对因子进行波动率调整 \\
\bottomrule
\end{longtable}
\end{readertable}

但是，\textbf{表达式数量的增长不等于独立信息量的增长。}

如果 Agent 只是不断组合相同的原始数据，那么它扩大的是 Expression Space 和 Hypothesis Space，而不是 IV-B 所定义的 Information Set。

这种区别对于解释后期性能退化尤为重要。

\hypertarget{part-4-63-0001-research-practice-factor-and-model-co-optimization}{%
\subsection{11. Research Practice --- Factor and Model Co-optimization}\label{part-4-63-0001-research-practice-factor-and-model-co-optimization}}

\hypertarget{part-4-63-0001-two-optimization-dimensions}{%
\subsubsection{11.1 Two Optimization Dimensions}\label{part-4-63-0001-two-optimization-dimensions}}

沿用前面的 Agent Version \(A_k\)，假设每个版本包含一个因子生成方法和一个预测模型。

定义：
\[f_k:\text{Raw Data}\longrightarrow X_{i,t}^{[k]},\]

\[h_k:X_{i,t}^{[k]}
\longrightarrow\hat r_{i,t+1}^{[k]}.\]

其中上标 \([k]\) 表示 Agent 版本，以区别此前用于标记单个因子的上标 \((j)\)。

那么整体预测流程为：
\[\boxed{
\hat r_{i,t+1}^{[k]}
=
h_k\circ f_k(\text{Raw Data}_{i,t})
}\]

Agent 的一次 Revision 可能改变 \(f_k\)，也可能改变 \(h_k\)，或者同时改变两者。

\hypertarget{part-4-63-0001-factormodel-interaction-experiment}{%
\subsubsection{11.2 Factor--Model Interaction Experiment}\label{part-4-63-0001-factormodel-interaction-experiment}}

假设 RSI 从 Baseline Version 进入 New Version，并同时修改了因子生成器和预测模型。

我们可以构建以下四组交叉实验。

Group A

\begin{verbatim}
    **Original Factors**

    Original Model
\end{verbatim}

\textbf{0.040}

OOS Mean IC

Group B

\begin{verbatim}
    **New Factors**

    Original Model
\end{verbatim}

\textbf{0.025}

OOS Mean IC

Group C

\begin{verbatim}
    **Original Factors**

    New Model
\end{verbatim}

\textbf{0.050}

OOS Mean IC

Group D

\begin{verbatim}
    **New Factors**

    New Model
\end{verbatim}

\textbf{0.060}

OOS Mean IC

假想实验结果。各组合均需要在相同数据协议下重新训练模型，而不是直接将旧权重用于不同输入特征。

这里，新因子单独使用时损害预测能力，但与新模型结合之后反而产生了更好的表现。

按照 III-B 中学习的 Factorial Experiment：
\[\begin{aligned}
\beta_{FM}
&=\mu_D-\mu_B-\mu_C+\mu_A\\
&=0.060-0.025-0.050+0.040\\
&=0.025.
\end{aligned}\]

这个结果描述了正向的 Factor--Model Interaction。

当然，真正判断其统计显著性还需要重复实验与不确定性估计。

这类实验很适合用于 RSI Revision Diagnosis，因为它能够区分性能变化来自新因子、新模型，还是二者之间的相互作用。

\hypertarget{part-4-63-0001-a-minimal-factor-evaluation-in-python}{%
\subsubsection{11.3 A Minimal Factor Evaluation in Python}\label{part-4-63-0001-a-minimal-factor-evaluation-in-python}}

以下代码给出了金融面板数据中计算每日 IC、Rank IC 和 ICIR 的基本方法。

\begin{verbatim}
import numpy as np
import pandas as pd


def evaluate_factor(
    df,
    factor_col="factor",
    label_col="future_return",
    date_col="date"
):
    """
    Each row represents one asset at one date.
    Labels must be correctly aligned in time.
    """

    records = []

    for date, group in df.groupby(date_col):
        data = group[
            [factor_col, label_col]
        ].dropna()

        if len(data) < 3:
            continue

        x = data[factor_col]
        y = data[label_col]

        if x.nunique() < 2 or y.nunique() < 2:
            continue

        records.append({
            "date": date,
            "ic": x.corr(y, method="pearson"),
            "rank_ic": x.corr(y, method="spearman"),
            "n_assets": len(data)
        })

    if not records:
        raise ValueError("No valid cross-sections")

    result = (
        pd.DataFrame(records)
        .set_index("date")
        .sort_index()
    )

    mean_ic = result["ic"].mean()
    std_ic = result["ic"].std(ddof=1)

    metrics = {
        "mean_ic": mean_ic,
        "mean_rank_ic": result["rank_ic"].mean(),
        "ic_std": std_ic,
        "icir": (
            mean_ic / std_ic
            if std_ic > 0
            else np.nan
        ),
        "positive_ic_rate": (
            result["ic"] > 0
        ).mean()
    }

    return result, metrics
\end{verbatim}

需要特别强调，计算这些指标本身并不困难，困难的是保证输入数据不存在 Look-ahead Bias、Survivorship Bias 或错误的标签对齐。

这段代码属于基础参考实现，尚未实际运行，也不是包含统计检验的完整 Factor Evaluation Framework。

\hypertarget{part-4-63-0001-research-hypothesis-factor-complexity-and-oos-decay}{%
\subsection{12. Research Hypothesis --- Factor Complexity and OOS Decay}\label{part-4-63-0001-research-hypothesis-factor-complexity-and-oos-decay}}

现在可以将本阶段知识连接到 X Program \#161 的核心研究问题。

假设我们定义：
\[p_k=\text{number of factors at iteration }k,\]

\[Q_k=\text{cumulative evaluation calls},\]

\[\Omega(f_k)=\text{factor expression complexity}.\]

那么在 Instrumentation 中，除了记录：
\[S_k^{\mathrm{IS}},
\qquad
S_k^{\mathrm{OOS}},\]

还可以同时记录：
\[p_k,\qquad Q_k,\qquad\Omega(f_k).\]

\bookgraphic{0.786}{part-4-63-0001-026-01}{}

\bookgraphic{0.786}{part-4-63-0001-027-01}{}

两张图均为假想数据，用于说明 Factor Complexity Accumulation 与 Generalization Decay 之间可能存在的关联，而非实际实验结果。

如果观察到因子数量和复杂度持续增长，而 OOS 性能逐渐下降，我们可以提出 Factor Complexity Accumulation 的候选解释。

但是，因子数量可能只是与退化同步变化，而非退化的真正原因。

因此，需要通过以下实验设计进行区分。

\begin{readertable}
\begin{longtable}[]{@{}
  >{\raggedright\arraybackslash}p{(\columnwidth - 4\tabcolsep) * \real{0.29}}
  >{\raggedright\arraybackslash}p{(\columnwidth - 4\tabcolsep) * \real{0.31}}
  >{\raggedright\arraybackslash}p{(\columnwidth - 4\tabcolsep) * \real{0.36}}@{}}
\toprule
Experiment & Controlled Change & Research Question \\
\midrule
\endhead
Factor Count Control & 限制保留因子数量 & 是否由因子数量膨胀导致退化？ \\
Expression Depth Control & 限制表达式树深度 & 是否由复杂的组合规则导致退化？ \\
Redundancy Filtering & 去除高相关或重复因子 & 是否由冗余信息导致退化？ \\
Model Ablation & 固定因子，仅修改模型 & 模型结构的贡献是什么？ \\
Factor Ablation & 固定模型类别，仅修改因子 & 因子生成的贡献是什么？ \\
Budget Matching & 匹配评估次数 & 是否仅仅来自搜索规模扩大？ \\
\bottomrule
\end{longtable}
\end{readertable}

特别是 Expression Tree Complexity，可以考虑使用节点数量、最大深度和运算符数量等指标进行刻画。

但它们只是可测量的复杂度代理，并不直接等于 Statistical Learning Theory 中的 Hypothesis Class Complexity。

\hypertarget{part-4-63-0001-literature-index}{%
\subsection{13. Literature Index}\label{part-4-63-0001-literature-index}}

本阶段推荐以下四篇正式发表的同行评审研究，已经核对其出版来源。

\textbf{P0 --- Empirical Asset Pricing via Machine Learning}

Gu, Kelly \& Xiu · 2020 · The Review of Financial Studies

理解现代金融机器学习、非线性预测及 OOS Performance Evaluation。

\href{https://doi.org/10.1093/rfs/hhaa009}{DOI: 10.1093/rfs/hhaa009}

\textbf{P0 --- \ldots{} and the Cross-Section of Expected Returns}

Harvey, Liu \& Zhu · 2016 · The Review of Financial Studies

理解 Factor Zoo、Multiple Testing 与金融因子发现中的统计偏差。

\href{https://academic.oup.com/rfs/article-abstract/29/1/5/1843824?login=false}{DOI: 10.1093/rfs/hhv059}

\textbf{P1 --- Taming the Factor Zoo: A Test of New Factors}

Feng, Giglio \& Xiu · 2020 · The Journal of Finance

理解如何评价新因子相对于大量既有因子的边际贡献。

\href{https://onlinelibrary.wiley.com/doi/full/10.1111/jofi.12883}{DOI: 10.1111/jofi.12883}

\textbf{P1 --- Common Risk Factors in the Returns on Stocks and Bonds}

Fama \& French · 1993 · Journal of Financial Economics

理解传统资产定价因子的定义，以及 Risk Exposure 与 Alpha 的关系。

\href{https://doi.org/10.1016/0304-405X\%2893\%2990023-5}{DOI: 10.1016/0304-405X(93)90023-5}

目前建议优先阅读 Harvey、Liu 和 Zhu 的研究，因为它与我们关注的自动化搜索、多重检验和虚假发现具有最直接的联系。

\hypertarget{part-4-63-0001-iv-c-conclusion-index}{%
\subsection{14. IV-C --- Conclusion Index}\label{part-4-63-0001-iv-c-conclusion-index}}

\hypertarget{part-4-63-0001-core-conclusions}{%
\subsubsection{14.1 Core Conclusions}\label{part-4-63-0001-core-conclusions}}

\begin{readertable}
\begin{longtable}[]{@{}
  >{\raggedright\arraybackslash}p{(\columnwidth - 4\tabcolsep) * \real{0.12}}
  >{\raggedright\arraybackslash}p{(\columnwidth - 4\tabcolsep) * \real{0.67}}
  >{\raggedright\arraybackslash}p{(\columnwidth - 4\tabcolsep) * \real{0.18}}@{}}
\toprule
ID & Key Conclusion & Index \\
\midrule
\endhead
C1 & Alpha 依赖于所采用的风险模型 & §2 \\
C2 & Risk Factor 与 Predictive Factor 不完全等价 & §3 \\
C3 & 横截面排序准确不等于绝对收益预测准确 & §4 \\
C4 & Panel Data 不一定包含大量独立样本 & §4.3 \\
C5 & IC 衡量横截面线性预测相关性 & §5.1 \\
C6 & Rank IC 衡量预测排序质量 & §5.2 \\
C7 & Rank IC 对严格单调递增变换保持不变 & §5.3 \\
C8 & ICIR 不等于统计显著性 & §5.4 \\
C9 & Pearson Correlation 与最优线性校准后的 \(R^2\) 存在数学关系 & §6 \\
C10 & Factor Zoo 与 Multiple Testing 密切相关 & §7 \\
C11 & Ridge 缓解共线性与估计不稳定 & §8.2--3 \\
C12 & Lasso 可以产生稀疏参数估计 & §8.4 \\
C13 & 非线性模型可以发现边际相关性遗漏的交互规律 & §9 \\
C14 & 自动化因子挖掘可以表示为代码或表达式空间搜索 & §10 \\
C15 & Factor Improvement 与 Model Improvement 可能存在交互作用 & §11 \\
C16 & 表达式复杂度增长不等于新增信息 & §10、§12 \\
\bottomrule
\end{longtable}
\end{readertable}

\hypertarget{part-4-63-0001-core-formula-sheet}{%
\subsubsection{14.2 Core Formula Sheet}\label{part-4-63-0001-core-formula-sheet}}

\textbf{Result 1 --- Multi-factor Model}

\[r_{i,t}-r_{f,t}
=
\alpha_i+
\sum_{j=1}^{p}\beta_{i,j}F_{j,t}
+\varepsilon_{i,t}.\]

\textbf{Result 2 --- Information Coefficient}

\[\mathrm{IC}_t
=
\operatorname{Corr}_i
(\hat r_{i,t+1},r_{i,t+1}).\]

\textbf{Result 3 --- Rank IC}

\[\mathrm{RankIC}_t
=
1-\frac{6\sum_i d_i^2}{N_t(N_t^2-1)}.\]

无并列名次时成立。

\textbf{Result 4 --- Maximum Predictive Correlation}

\[\boxed{
|\operatorname{Corr}(r_{t+1},h(X_t))|
\leq
\sqrt{\frac{\mathrm{SNR}}{1+\mathrm{SNR}}}
}\]

\textbf{Result 5 --- Ridge Estimator}

\[\boxed{
\hat{\boldsymbol\beta}_{\mathrm{Ridge}}
=
(X^\top X+\lambda I_p)^{-1}X^\top\boldsymbol y
}\]

\textbf{Result 6 --- Lasso Soft Thresholding}

\[\hat\beta_j^{\mathrm{Lasso}}
=
\operatorname{sign}(z_j)
\left(|z_j|-\frac{\lambda}{2}\right)_+.\]

在本节目标函数定义及正交设计条件下成立。

\textbf{Result 7 --- Factor--Model Composition}

\[\hat r_{i,t+1}^{[k]}
=
h_k\circ f_k(\text{Raw Data}_{i,t}).\]

\hypertarget{part-4-63-0001-proof-index}{%
\subsubsection{14.3 Proof Index}\label{part-4-63-0001-proof-index}}

\begin{readertable}
\begin{longtable}[]{@{}
  >{\raggedright\arraybackslash}p{(\columnwidth - 4\tabcolsep) * \real{0.15}}
  >{\raggedright\arraybackslash}p{(\columnwidth - 4\tabcolsep) * \real{0.59}}
  >{\raggedright\arraybackslash}p{(\columnwidth - 4\tabcolsep) * \real{0.21}}@{}}
\toprule
Proof & Mathematical Result & Location \\
\midrule
\endhead
P1 & 共同风险因子导致资产相关性 & §4.3 \\
P2 & Rank Correlation 的单调变换不变性 & §5.3 \\
P3 & 最优线性校准与 Pearson Correlation & §6.2 \\
P4 & SNR 与预测相关性的理论上界 & §6.3 \\
P5 & Ridge Closed-form Solution & §8.2 \\
P6 & Ridge 的特征值解释 & §8.3 \\
P7 & 正交条件下的 Lasso Soft Thresholding & §8.4 \\
P8 & 非线性交互与零边际相关性 & §9 \\
\bottomrule
\end{longtable}
\end{readertable}

\hypertarget{part-4-63-0001-the-most-important-research-insight}{%
\subsection{15. The Most Important Research Insight}\label{part-4-63-0001-the-most-important-research-insight}}

到这里，我们已经能够更加具体地理解 Financial RSI 可能在优化什么。

\textbf{Fixed Raw Financial Information}

历史价格、交易量、企业基本面

↓

\textbf{Factor Optimization}

发现更有效的信息表示

\(f_k\rightarrow f_{k+1}\)

\textbf{Model Optimization}

改善预测函数

\(h_k\rightarrow h_{k+1}\)

↓ \textbf{Increasing Search \& Representation Capability} ↓

Genuine Improvement

发现有效信号

识别非线性交互

降低估计误差

Potential Overfitting

发现虚假因子

积累冗余特征

过度适应开发环境

↓ \textbf{Better, Then Worse}

取决于两种机制在递归演化过程中的相对影响

我们现在可以进一步提出一个值得研究的观点：

\textbf{Financial RSI 的困难之一，是 Agent 可以近乎无限地生成新的金融表达式，但这些表达式所依赖的真实信息可能始终有限。}

优化器的能力可以继续增长，搜索空间可以继续扩大，然而给定信息集合下的 Bayes Optimal Performance 并不会因此自动提高。

这仍然是我们根据统计理论提出的候选机制。课题组是否采用因子挖掘、横截面收益预测或其他金融任务，目前尚未确认。

\hypertarget{part-4-63-0001-next-iv-d-returns-sharpe-ratio-performance-evaluation}{%
\subsection{Next --- IV-D: Returns, Sharpe Ratio \& Performance Evaluation}\label{part-4-63-0001-next-iv-d-returns-sharpe-ratio-performance-evaluation}}

IV-A --- Financial Time Series \& DGP

IV-B --- Low Signal-to-Noise Ratio \& Predictability

IV-C --- Financial Factors \& ML Models

IV-D --- Returns, Sharpe Ratio \& Evaluation

IV-E --- Backtest Overfitting, PBO \& DSR

IV-F --- Market Regime \& Robustness

下一阶段我们将研究金融评价指标本身。这是理解 Backtest Overfitting 之前不可缺少的基础。

\begin{readertable}
\begin{longtable}[]{@{}
  >{\raggedright\arraybackslash}p{(\columnwidth - 2\tabcolsep) * \real{0.46}}
  >{\raggedright\arraybackslash}p{(\columnwidth - 2\tabcolsep) * \real{0.50}}@{}}
\toprule
Topic & Core Question \\
\midrule
\endhead
Portfolio Return & 如何将预测结果转化为组合收益？ \\
Expected Return \& Risk & 如何同时衡量收益和风险？ \\
Sharpe Ratio & 为什么不能只比较累计收益？ \\
Sharpe Ratio Estimation & Sharpe Ratio 本身具有怎样的统计误差？ \\
Annualization & 为什么乘以 \(\sqrt{252}\) 需要特定条件？ \\
Maximum Drawdown & 如何衡量最严重的历史回撤？ \\
Turnover \& Transaction Costs & 为什么预测能力不一定等于经济价值？ \\
Prediction vs. Economic Performance & 为什么高 IC 不必然带来高 Sharpe？ \\
\bottomrule
\end{longtable}
\end{readertable}

尤其值得关注的是，Sharpe Ratio 并非一个固定且精确的真实性能指标，而是一个存在抽样误差的统计估计量：
\[\widehat{\mathrm{SR}}
=
\frac{\hat\mu}{\hat\sigma}.\]

当 RSI Agent 不断以经验 Sharpe Ratio 作为优化目标时，它实际上是在反复优化一个由随机变量之比构成的统计量。

\textbf{IV-D 将研究这个统计量的性质；IV-E 则将进一步研究大规模搜索如何使它产生系统性高估，以及 PBO、Probabilistic Sharpe Ratio 和 Deflated Sharpe Ratio 如何处理这一问题。}
